%% Submitted to: Mathematics (MDPI)
%% Class: Definitions/mdpi  (download from https://www.mdpi.com/authors/latex)

\documentclass[journal,article,submit,pdftex,moreauthors]{Definitions/mdpi}

% ---------------------------------------------------------------
% Journal metadata  (fill in after acceptance)
% ---------------------------------------------------------------
\firstpage{1}
\makeatletter
\setcounter{page}{\@firstpage}
\makeatother
\pubvolume{1}
\issuenum{1}
\articlenumber{0}
\pubyear{2026}
\copyrightyear{2026}
\datereceived{}   % TODO
\daterevised{}    % TODO
\dateaccepted{}   % TODO
\datepublished{}  % TODO
% \hreflink{https://doi.org/}  % TODO: add DOI

% ---------------------------------------------------------------
% Extra packages
% (mdpi.cls already loads: amsmath, graphicx, hyperref, natbib,
%  geometry, color, xcolor, float, caption, etoolbox, and others.
%  Only add what the class does NOT provide.)
% ---------------------------------------------------------------
\usepackage{amssymb}
\usepackage{amsthm}
\usepackage{booktabs}       % \toprule / \midrule / \bottomrule
\usepackage{multirow}
\usepackage{adjustbox}
\usepackage{array}
\usepackage{algorithm}
\usepackage{algpseudocode}
\usepackage{pifont}
\newcommand{\cmark}{\ding{51}}
\newcommand{\xmark}{\ding{55}}

% ---------------------------------------------------------------
% Theorem environments
% (mdpi.cls defines Theorem, Lemma, Corollary, Proposition,
%  Remark, Definition, Example, Exercise, Hypothesis, Notation,
%  Problem, Property, Proof -- reuse those; add only Claim.)
% ---------------------------------------------------------------
\newtheorem{claim}{Claim}[section]

% ---------------------------------------------------------------
% ORCID
% mdpi.cls provides \orcidA{}, \orcidB{}, \orcidC{} etc.
% ---------------------------------------------------------------
\newcommand{\orcidauthorA}{0000-0001-6384-8840}   % Elad Shoham
\newcommand{\orcidauthorB}{0009-0005-9746-0058}   % Havana Rika
\newcommand{\orcidauthorC}{0000-0002-3895-688X}   % Dan Vilenchik

% ---------------------------------------------------------------
% Title, authors, affiliations
% ---------------------------------------------------------------
\Title{The Missing Signal: Gradient Injection Recovers $O(n^2)$ Residual Decoding in GNNs}

% Short title for running head
% \TitleCitation{Gradient Injection Recovers $O(n^2)$ Residual Decoding in GNNs}

% Authors: superscript number = affiliation; * = corresponding; \orcidX{} = icon
\Author{
  Elad Shoham $^{1}$\orcidA{},
  Havana Rika $^{2}$\orcidB{}
  and Dan Vilenchik $^{1,*}$\orcidC{}
}

% Plain-text author list for metadata
\AuthorNames{Elad Shoham, Havana Rika and Dan Vilenchik}

% Citation-style author list
% \AuthorCitation{Shoham, E.; Rika, H.; Vilenchik, D.}

% Affiliations (one per line, $^{N}$ prefix)
\address{
  $^{1}$ \quad Department of Software and Information Systems Engineering,
  Ben-Gurion University of the Negev,
  Beer-Sheva 8410501, Israel;
  shellad@post.bgu.ac.il (E.S.)\\
  $^{2}$ \quad The Academic College of Tel Aviv-Yaffo,
  Tel Aviv 61083, Israel;
  havanari@mta.ac.il (H.R.)
}

% Corresponding author
\corres{Correspondence: shellad@post.bgu.ac.il; vilenchi@bgu.ac.il}

% ---------------------------------------------------------------
% Abstract  (placed BEFORE \begin{document} in mdpi.cls)
% ---------------------------------------------------------------
\abstract{%
  Residual decoding algorithms score elements, prune one, and update state
  iteratively. Classical algorithms exploit additive statistics such as degree to
  perform each update in $O(n)$, achieving $O(n^2)$ total cost. Graph Neural
  Networks (GNNs) cannot replicate this efficiency: nonlinear activations such as
  $\tanh$ or sigmoid make each unit of score change unequal---the gap between
  $\tanh(6)$ and $\tanh(5)$ differs fundamentally from that between $\tanh(2)$ and
  $\tanh(1)$---so no fixed additive correction exists after a deletion. This forces
  a full network rerun after every deletion, inflating cost to $O(n^3)$.
  We show that the objective gradient provides the missing signal and validate this
  on \textsc{Max-Clique} using planted-clique instances: after each vertex deletion,
  each surviving gradient coordinate changes by a local, maintainable quantity,
  updatable in $O(n)$ per step. Supplying this live gradient to a lightweight
  learned updater recovers $O(n^2)$ decoding, matching or exceeding the $O(n^3)$
  rerun baseline across all difficulty regimes.
}

% Keywords (3--10; MDPI Mathematics style: lowercase, semicolon-separated)
\keyword{%
  graph neural networks;
  combinatorial optimization;
  Max-Clique;
  concept learning;
  unsupervised
}

% \MSC{}   % TODO: add MSC 2020 codes, e.g. 68R10; 90C27; 68T07

% ================================================================
\begin{document}
\maketitle
% ================================================================

% ================================================================
\section{Introduction}
% ================================================================

A wide class of combinatorial algorithms proceeds by a \emph{residual loop}: score
the current elements, prune or select one, update the problem state, and repeat.
This loop underlies greedy peeling~\citep{feige2010}, spectral
rounding~\citep{alon1998}, branch-and-bound solvers, and local greedy methods more
broadly. Its tractability rests on the \emph{additivity} of the maintained
statistics. Degree is the simplest example: when a vertex is removed, each
survivor's degree changes by at most one, depending only on local adjacency. This
local, fixed correction makes the full update $O(n)$, giving $O(n^2)$ total
decoding on dense graphs.


We focus on \textsc{Max-Clique}: given a graph, find the largest subset of vertices
that are all pairwise adjacent. The decision version is NP-complete and the
optimization problem is NP-hard, making clique recovery a notorious testbed for
algorithmic heuristics. Classical decoders attack it with greedy top-$k$ selection,
local-degree refinement, spectral rounding, and related residual scoring rules. In
the planted-clique setting used in our experiments, the clique is hidden inside a
random background graph and the decoder must recover it. This makes
\textsc{Max-Clique} a natural place to ask whether learned scorers can replace these
hand-designed residual heuristics without losing their efficient update structure.


GNNs have emerged as the natural candidate for replacing hand-crafted statistics
with learned scoring rules~\citep{karalias2020,min2022,garmendia2024}. Trained
end-to-end on problem instances, they produce vertex rankings for \textsc{Max-Clique},
\textsc{Max-Cut}, \textsc{Max-Independent-Set}, and related problems without
domain-specific feature engineering. Recent interpretability work shows that GNNs
trained on combinatorial tasks tend to rediscover classical algorithmic
heuristics~\citep{cappart2023}, and for \textsc{Max-Clique} specifically,
degree-based vertex ranking emerges consistently across architectures and training
regimes~\citep{shoham2026ltr}. This raises an optimistic picture: a GNN that has
internalized degree-based scoring should, in principle, slot directly into the
residual loop in place of the classical scorer.


This hope runs into a fundamental incompatibility. Classical residual algorithms work
because their maintained statistics update additively: removing a vertex changes each
survivor's degree by exactly one, so a single fixed correction suffices. GNN scores
cannot support this. Nonlinear activations such as $\tanh$ or sigmoid make each unit
of score change unequal---the gap between $\tanh(6)$ and $\tanh(5)$ differs from
that between $\tanh(2)$ and $\tanh(1)$---so no fixed additive correction exists
after a deletion. The problem compounds with scale: as instance size grows, score
gaps between adjacent candidates vanish, so even a small approximation error
eventually inverts the removal ordering. Once an inversion occurs, the approximate
and exact decoders evolve on structurally different residual graphs, and subsequent
errors grow without bound. The only safe remedy is to rerun the GNN after every
deletion---$O(n^2)$ per step, $O(n^3)$ total.

This leads us to our central question: \textbf{\emph{Can a GNN be designed to
    support efficient incremental decoding, recovering the $O(n^2)$ complexity of
    classical residual algorithms without giving up the expressive power of neural
scoring?}}


The obstruction above does not rule out efficient decoding; it rules out efficient
decoding from the \emph{raw GNN score} alone. The objective gradient provides
exactly the missing signal: after each vertex deletion, each surviving gradient
coordinate changes by a local, maintainable quantity determined entirely by the
interaction between the deleted vertex and each survivor. For the clique objective
this correction can be maintained in $O(n)$ per deletion by tracking three
sufficient statistics, giving $O(n^2)$ total cost after the initial GNN pass.
Supplying this live cached gradient to a learned score updater $U_\theta$ gives it
access to the residual correction directly, without requiring any inversion of the
underlying nonlinear activations.

\subsection{Contributions}

We identify a fundamental incompatibility between nonlinear activations and additive
residual updates: score gaps between adjacent candidates vanish with instance size,
so any fixed approximation error eventually inverts the removal ordering, after which
errors compound without bound. The objective gradient provides the missing signal,
as its per-deletion change is local and maintainable in $O(n)$ per step by tracking
three sufficient statistics, recovering $O(n^2)$ total decoding cost after the
initial GNN pass.

On the empirical side, we introduce U-LPR (Updatable Least Probable Removal), a
decoder that runs the GNN once and updates scores incrementally via a lightweight
learned updater $U_\theta$. We evaluate two variants: SGNN+U, which provides
$U_\theta$ with no gradient information, and SGNN+GU, which injects the cached
objective gradient into $U_\theta$ at each deletion step. We also introduce SIT
(Self-Imitation Training), a training regime that supervises $U_\theta$ by imitating
a rerun-based teacher. Both are described in full in Section~\ref{sec:method}.
Our findings:
\begin{itemize}
  \item \textbf{Recovery quality and running time.} SGNN+GU under SIT training
    achieves perfect approximation and exact recovery on Easy and Medium regimes
    at $O(n^2)$ cost, matching the $O(n^3)$ rerun baseline with a $4\times$
    wall-clock speedup.
  \item \textbf{Necessity of the gradient channel.} Removing the gradient from
    $U_\theta$ collapses U-LPR performance even under SIT training, confirming
    that gradient injection is a necessary architectural condition, not merely a
    helpful signal. Without it, SIT training makes performance strictly worse than
    objective training alone.
  \item \textbf{Learned residual concept.} Interpretability analysis reveals that
    the successful updater learns a coarse low-degree predicate: it avoids deleting
    vertices above the 15th percentile of residual degree among alive nodes, rather
    than sorting vertices by exact degree. The gradient channel is what makes this
    concept actionable during the residual loop.
\end{itemize}

Figure~\ref{fig:decoder_overview} summarizes the architectural distinction between
SGNN, SGNN+U, and SGNN+GU.

\begin{figure}[H]
  \centering
  \includegraphics[width=0.92\textwidth]{figures/decoder_figure.png}
  \caption{Overview of the three neural decoder variants. SGNN uses the backbone
    scores directly, SGNN+U adds an incremental score updater without gradient
    information, and SGNN+GU supplies the updater with a cached residual gradient.
    The gradient channel is the mechanism that makes the learned updater compatible
  with residual decoding.}
  \label{fig:decoder_overview}
\end{figure}

\subsection{Paper Organization}

Section~\ref{sec:related} reviews related work. Section~\ref{sec:method} develops
the methodology and introduces U-LPR and SIT. Section~\ref{sec:theory} presents
the analysis of why nonlinear scores fail and how the gradient resolves this.
Section~\ref{sec:experiments} reports the empirical evaluation.
Section~\ref{sec:concepts} analyzes the learned residual concept.
Section~\ref{sec:discussion} concludes with limitations and future directions.

% ================================================================
\section{Related Work}
\label{sec:related}
% ================================================================

Our backbone is a message-passing neural network (MPNN), which maintains a hidden
state $h_v^{(\ell)} \in \mathbb{R}^d$ updated over $L$ layers as
\[
  h_v^{(\ell+1)}
  =
  \sigma\!\left(
    \mathbf{M}\!\left[
      h_v^{(\ell)},\,
      \mathrm{Agg}_{u\in\mathcal{N}(v)} h_u^{(\ell)}
    \right]
  \right),
\]
with a per-vertex score $z_v = \mathbf{W}_{\mathrm{out}} h_v^{(L)}$ read off at the
end~\citep{xu2019}, optionally passed through a final nonlinearity such as sigmoid
to produce vertex probabilities $p_v \in (0,1)$. Their expressive power is bounded
by the 1-WL test~\citep{xu2019,morris2019}, though our results concern an orthogonal
limitation---\emph{update compatibility}---that persists even beyond the WL bound.


\citet{cappart2023} argue that GNNs on NP-hard problems behave more like
polynomial-time heuristics than exact solvers, and \citet{shoham2026ltr} give a
concrete \textsc{Max-Clique} example: across architectures and datasets, learned
representations consistently recover degree-based vertex rankings. The difficulty is
that a useful residual heuristic is not only a good ranking rule; it is also a rule
whose state can be updated cheaply after each deletion. Our work shows that this
second property does not automatically transfer from degree to a nonlinear GNN score.


\citet{karalias2020} introduced an unsupervised framework deriving GNN loss directly
from the problem objective; \citet{min2022} extended this with the scattering-based
architecture we build on. Similar unsupervised objectives have been applied to
\textsc{Max-Cut}, \textsc{Max-Independent-Set}, and
\textsc{TSP}~\citep{schuetz2022,garmendia2024}. These works demonstrate strong
combinatorial heuristics, but the standard decoding pattern still reruns or reuses
the network in a residual loop---precisely where nonlinear activations make score
changes unequal across the activation range, so no fixed additive correction exists.


\citet{yau2023} (OptGNN) proved that polynomial-sized MPNNs can implement SDP
gradient descent for Max-CSPs, obtaining UGC-optimal guarantees in theory and
outperforming SDP solvers in practice. QRF-GNN~\citep{gaile2024} feeds intermediate
GNN predictions back as dynamic node features to improve convergence within a single
inference pass, and \citet{hadou2024} unroll a dual ascent algorithm into coupled
networks so each layer computes a meaningful algorithmic step. Our work is
complementary: rather than aligning message passing with a full relaxation
trajectory, we align the residual decoder with the local objective change caused by
each deletion.


This connects our SIT regime to imitation learning in combinatorial
optimization~\citep{liu2021}. The novelty here is self-imitation: teacher and
student share weights, and the student must predict the rerun-based removal sequence
from incrementally updated scores using the cached gradient.

% ================================================================
\section{Methodology}
\label{sec:method}
% ================================================================

\subsection{Problem Setting}
\label{sec:problem}

We focus on \textsc{Max-Clique}: given a graph $G=(V,E)$, find the largest subset
$C \subseteq V$ such that all pairs in $C$ are adjacent. We evaluate on the
planted-clique distribution $G(n,\tfrac{1}{2},k)$: an Erd\H{o}s--R\'enyi random
graph on $n$ vertices where each edge is included independently with probability
$1/2$, with a clique $C^\star$ of size $k$ planted by adding all missing edges among
$k$ uniformly chosen vertices. The goal is to recover a large valid clique, ideally
$C^\star$ itself.

\subsubsection{Decoders}

A natural family of decoders for this problem is the \emph{residual loop}: score all
vertices, remove the lowest-scoring one, and repeat until the remaining set forms a
clique, followed by an add-back pass that restores any removed vertex compatible with
the found clique. We evaluate five decoders:
\begin{itemize}
  \item \textbf{Top-$k$}: a non-iterative baseline that runs the GNN once, ranks
    vertices by score, and greedily adds them to the solution if adjacent to all
    previously selected vertices. $O(n^2)$.
  \item \textbf{LDR/LPR} (Least Degree/Probable Removal): the classical
    instantiation scores vertices by their current residual degree; LPR replaces
    degree with a learned probability $p_v$ produced by a GNN, retaining the same
    residual loop structure with scores fixed after the initial pass. Both are
    $O(n^2)$.
  \item \textbf{Rerun}: reruns the full GNN after every deletion to produce fresh
    scores. The only $O(n^3)$ decoder, serving as the adaptive upper bound.
  \item \textbf{U-LPR} (Updatable Least Probable Removal): runs the GNN once and
    applies a learned MLP updater $U_\theta$ after each deletion to incrementally
    correct the scores (Algorithm~\ref{alg:gub}). $O(n^2)$.
\end{itemize}

\subsubsection{Datasets and Difficulty Regimes}

We evaluate on synthetic planted-clique instances $G(n,\tfrac{1}{2},k)$ with
$n=1000$ vertices. Each instance is generated by first sampling $k$ uniformly from
the relevant regime interval, then choosing $k$ vertices uniformly at random and
adding all missing edges among them; all remaining edges are sampled independently
with probability $1/2$.

Following~\citet{shoham2026ltr}, the Easy/Medium/Hard labels are derived from the
algorithmic thresholds for planted clique recovery and then calibrated at the fixed
finite graph size. Asymptotically, larger planted cliques give the planted vertices
a visible degree boost: when $k=\Omega(\sqrt{n\log n})$, simple degree-based
ranking succeeds; when $k=\Omega(\sqrt n)$ but below the degree-ranking threshold,
recovery requires stronger residual or spectral methods; and around the
$O(\sqrt n)$ scale, the planted clique hardness conjecture predicts that recovery
is computationally intractable.

For a fixed $n$, the finite thresholds are obtained empirically as follows. Let
$k_1$ be the smallest clique size for which the Top-$k$ degree heuristic recovers
the planted clique at least $90\%$ of the time. Lower $k$ until degree-based LPR
falls below $90\%$ success, and denote that cutoff by $k_2$. Finally, set
$k_3=\lceil 2\log_2 n\rceil$, the typical largest-clique scale in the unplanted
background graph~\citep{bollobas1985}. The full calibrated regions are Easy
$[k_1,n]$, Medium $[k_2+1,k_1-1]$, and Hard $[k_3,k_2]$.

For $n=1000$, this calibration gives $k_1=62$, $k_2=35$, and $k_3=20$. The reported
experiments therefore use the following sampled ranges:
\begin{itemize}
  \item \textbf{Easy} ($k\in[62,100]$): a bounded subset of the full easy region
    $k\ge 62$, where degree-based ranking already succeeds~\citep{kucera1995}.
  \item \textbf{Medium} ($k\in[36,61]$): below the Top-$k$ degree threshold but
    still recoverable by stronger residual or spectral methods such
    as degree-based LPR and AKS~\citep{alon1998,feige2010}.
  \item \textbf{Hard} ($k\in[20,35]$): near the background clique scale and below
    the empirical residual-method threshold, where no polynomial-time method is
    known to recover the planted clique reliably~\citep{manurangsi2021}.
\end{itemize}

\subsection{Model Architecture}
\label{sec:architectures}

We evaluate three model families: SGNN, SGNN+U, and SGNN+GU. All share the same
message-passing backbone with $L=4$ steps, hidden dimension $64$, and
$\sigma=\tanh$. The input feature is a scalar liveness indicator; no degree,
clique-membership, or structural feature is provided. The backbone aggregates and
updates as
\[
  q_v^{(\ell)} = \frac{1}{n-1}\sum_{u\in\mathcal{N}(v)} h_u^{(\ell)},
  \qquad
  h_v^{(\ell+1)} = \sigma\!\left(M_{\mathrm{upd}}\!\left[h_v^{(\ell)},
  q_v^{(\ell)}\right]\right),
\]
\[
  s_v = W_{\mathrm{out}} h_v^{(L)} + b_{\mathrm{out}},
  \qquad
  p_v = \operatorname{sigmoid}(s_v),
\]
where $M_{\mathrm{upd}}$ is a two-layer MLP and $p_v \in (0,1)$ is the vertex
probability used by the clique objective. SGNN uses this backbone alone. SGNN+U and
SGNN+GU add a lightweight MLP updater $U_\theta$ that corrects surviving scores
after each deletion of $v^*$:
\[
  s_i \leftarrow s_i + U_\theta(s_i,\, s_{v^*},\, A_{iv^*},\, \gamma_i,\,
  \gamma_{v^*}),
\]
where $\gamma_i, \gamma_{v^*}$ are gradient channels that differ between variants.
For SGNN+U, $\gamma_i = \gamma_{v^*} = 0$. For SGNN+GU,
\[
  \gamma_i = g_i = \frac{\partial \mathcal{L}_{\mathrm{clique}}}{\partial p_i},
  \qquad
  \gamma_{v^*} = g_{v^*} = \frac{\partial \mathcal{L}_{\mathrm{clique}}}{\partial
  p_{v^*}},
\]
maintained via cached sufficient statistics updated in $O(n)$ per deletion
(Section~\ref{sec:cached_maintenance}).

\subsection{Cached Gradient Maintenance}
\label{sec:cached_maintenance}

Recomputing the full gradient after every deletion costs $O(n^2)$ per step and
$O(n^3)$ total. For our clique objective, this is unnecessary. The gradient at
each surviving vertex is a linear function of two cached sums: one aggregating
probabilities over neighbors, and one over non-neighbors. When a vertex $v$ is
deleted, its contribution to each sum is simply subtracted---an $O(n)$ operation
per deletion rather than a full $O(n^2)$ recomputation. The cache is initialized
in $O(n^2)$ alongside the initial GNN pass, so the total decoding cost remains
$O(n^2)$.

\subsection{Training}
\label{sec:training}

All models are trained on Medium-difficulty instances and evaluated across all three
regimes. Each run uses 64 epochs, with 1000 newly sampled instances per epoch,
minibatch size 32, and no overlap with the held-out evaluation instances. Training
on Easy or Hard instances yields identical results, confirming that the findings are
not sensitive to the choice of training regime. The unsupervised clique objective is
\begin{equation}
  \mathcal{L}_{\mathrm{clique}}(p; A, k)
  =
  -p^\top \mathbf{A} p
  +
  p^\top \mathbf{B} p,
  \label{eq:clique_loss}
\end{equation}
where $p_i = \sigma(s_i)$ and $\mathbf{B}$ is the complement adjacency matrix.

\subsubsection{Objective Training (OBJ)}

$\mathcal{L}_{\mathrm{train}} = \mathcal{L}_{\mathrm{clique}}$. $U_\theta$ is
supervised only through the objective, with no direct signal on the removal
sequence.

\subsubsection{Self-Imitation Training (SIT)}

SIT adds an imitation loss that directly supervises the incremental decoder. At each
step $t$, a rerun-based teacher reruns the GNN on the current residual graph $G_t$
and selects the lowest-scoring vertex $v^t_{\mathrm{teach}}$. A student---running
the GNN only once and updating scores incrementally via $U_\theta$---is penalised
for disagreeing:
\begin{equation}
  \mathcal{L}_{\mathrm{SIT}} = \frac{1}{T} \sum_{t=0}^{T-1}
  \mathrm{CE}_{\mathrm{softmax}}\!\left(-\tilde{s}^{(t)}_{V(G_t)},\;
  v^t_{\mathrm{teach}}\right),
  \qquad
  \mathcal{L}_{\mathrm{train}} = \mathcal{L}_{\mathrm{clique}} +
  \mathcal{L}_{\mathrm{SIT}}.
  \label{eq:sit_loss}
\end{equation}
For SGNN+U, $U_\theta$ receives no gradient signal and is trained with
$\gamma_i = \gamma_{v^*} = 0$. For SGNN+GU, the cached gradient is updated in
$O(n)$ after each teacher removal before being passed to $U_\theta$. SIT uses
$T=300$ teacher removals per instance; the teacher cost is paid at train time only.

% ================================================================
\section{Claims and Proofs}
\label{sec:theory}
% ================================================================

We formalize the three arguments sketched in the introduction.

\subsection{Why Nonlinear Scores Fail}

\begin{claim}
  Exact additive updates require the score encoding to be affine in the maintained
  statistic. Nonlinear activations such as $\tanh$ or sigmoid are incompatible with
  this: the gap between $\tanh(6)$ and $\tanh(5)$ differs from the gap between
  $\tanh(2)$ and $\tanh(1)$, so no fixed additive correction exists after a
  deletion.
\end{claim}
\begin{proof}
  If $\varphi(d-1) = \varphi(d) - c$ for a fixed $c$ and all $d$, then every first
  discrete difference of $\varphi$ is constant, forcing $\varphi(d) = cd + b$. Any
  nonlinear activation $\sigma$ satisfies $\sigma(a-b) \neq \sigma(a) - \sigma(b)$
  in general, so no choice of parameters can make the post-deletion score equal the
  pre-deletion score minus a fixed constant.
\end{proof}

\begin{claim}
  Any fixed per-step approximation error $\varepsilon > 0$ eventually inverts the
  removal ordering. Once an inversion occurs, the approximate and exact decoders
  diverge onto structurally different residual graphs, after which errors compound
  without bound.
\end{claim}
\begin{proof}
  Let $\gamma_n$ be the minimum score gap. After $r = \lceil\gamma_n/\varepsilon\rceil$
  steps the accumulated error $r\varepsilon$ exceeds $\gamma_n$, inverting a pairwise
  ordering. After the inversion the two decoders remove different vertices; every
  subsequent score is evaluated on a different residual graph, so the discrepancy is
  no longer a perturbation of the original error but an uncontrolled structural
  difference.
\end{proof}

\subsection{The Gradient as the Missing Signal}

\begin{claim}
  For any locally decomposable objective, deleting a vertex removes exactly its
  local contribution from each surviving gradient coordinate. For quadratic
  objectives this correction is linear in the deleted vertex's cached probability,
  and is directly representable in the first linear layer of $U_\theta$ before any
  nonlinearity fires.
\end{claim}
\begin{proof}
  Differentiating a locally decomposable objective, the gradient at vertex $i$
  involves only terms that include $i$. Deleting vertex $v$ removes exactly the
  pairwise term between $i$ and $v$, leaving all other terms unchanged. For a
  quadratic objective the removed term is linear in $p_v$, which is available as an
  input to $U_\theta$. The first affine layer of $U_\theta$ can therefore represent
  the correction exactly before any activation is applied. Without the gradient,
  recovering this correction from the nonlinear score $\sigma(z_v)$ requires
  inverting $\sigma$, which the first claim shows is impossible with a fixed
  additive rule.
\end{proof}

\subsection{Cached Maintenance Recovers $O(n^2)$}

\begin{claim}
  The gradient cache can be initialized in $O(n^2)$ and updated in $O(n)$ per
  deletion, giving $O(n^2)$ total decoding cost after the initial GNN pass.
\end{claim}
\begin{proof}
  The gradient depends on two cached sums over neighbors and non-neighbors, each
  initialized by a matrix-vector product in $O(n^2)$. Deleting vertex $v$ subtracts
  its contribution from each sum in $O(n)$. Over at most $n$ deletions the total
  update cost is $O(n^2)$, matching the initial pass.
\end{proof}

% ================================================================
\section{Experiments}
\label{sec:experiments}
% ================================================================

\subsection{Setup}

We use the architectures from Section~\ref{sec:architectures} and the training
protocol from Section~\ref{sec:training}. All experiments use planted-clique
instances with $n=1000$; evaluation set sizes are reported in the table captions.
Reported values are mean $\pm$ standard deviation over the corresponding held-out
instances. Hardware: AMD Ryzen 9 5950X, 64\,GB RAM, NVIDIA RTX 3080; Python 3.11,
PyTorch 2.2.

\subsection{Main Results}
\label{sec:main_results}

Table~\ref{tab:main} reports planted clique recovery across all model and decoder
combinations.

\begin{table}[H]
  \caption{Planted clique recovery on $n=1000$ across Easy, Medium, and Hard
    regimes, reported as mean$\,\pm\,$std over 10 held-out instances per regime;
    entries without $\pm$ are exactly $0.00$ or $1.00$ across all instances.
    \textbf{Bold}: best $O(n^2)$ result per regime. SGNN has no updater and cannot
  run U-LPR.}
  \label{tab:main}
  \resizebox{\textwidth}{!}{%
    \begin{tabular}{lllccccccccc}
      \toprule
      & & & \multicolumn{3}{c}{Easy} & \multicolumn{3}{c}{Medium}
      & \multicolumn{3}{c}{Hard} \\
      \cmidrule(lr){4-6}\cmidrule(lr){7-9}\cmidrule(lr){10-12}
      Model & Decoder & Runtime & Approx & Overlap & Exact
      & Approx & Overlap & Exact
      & Approx & Overlap & Exact \\
      \midrule
      \multirow{3}{*}{SGNN + OBJ}
      & Top-$k$   & $O(n^2)$ & $1.00$ & $1.00$ & $1.00$
      & $0.75{\pm}0.43$ & $0.74{\pm}0.45$ & $0.67{\pm}0.58$
      & $0.46{\pm}0.12$ & $0.08{\pm}0.07$ & $0.00$ \\
      & Rerun     & $O(n^3)$ & $1.00$ & $1.00$ & $1.00$
      & $1.00$ & $1.00$ & $1.00$
      & $0.49{\pm}0.09$ & $0.00$ & $0.00$ \\
      & LPR       & $O(n^2)$ & $1.00$ & $1.00$ & $1.00$
      & $0.75{\pm}0.43$ & $0.74{\pm}0.45$ & $0.67{\pm}0.58$
      & $0.46{\pm}0.12$ & $0.08{\pm}0.07$ & $0.00$ \\
      \midrule
      \multirow{4}{*}{SGNN+U + OBJ}
      & Top-$k$   & $O(n^2)$ & $1.00$ & $1.00$ & $1.00$
      & $0.62{\pm}0.36$ & $0.58{\pm}0.40$ & $0.33{\pm}0.58$
      & $0.37{\pm}0.06$ & $0.08{\pm}0.07$ & $0.00$ \\
      & Rerun     & $O(n^3)$ & $1.00$ & $1.00$ & $1.00$
      & $1.00$ & $1.00$ & $1.00$
      & $0.63{\pm}0.33$ & $0.36{\pm}0.56$ & $0.33{\pm}0.58$ \\
      & LPR       & $O(n^2)$ & $1.00$ & $1.00$ & $1.00$
      & $0.62{\pm}0.36$ & $0.58{\pm}0.40$ & $0.33{\pm}0.58$
      & $0.37{\pm}0.06$ & $0.08{\pm}0.07$ & $0.00$ \\
      & U-LPR     & $O(n^2)$ & $0.13{\pm}0.01$ & $0.00$ & $0.00$
      & $0.17{\pm}0.01$ & $0.00$ & $0.00$
      & $0.34{\pm}0.09$ & $0.00$ & $0.00$ \\
      \midrule
      \multirow{4}{*}{SGNN+GU + OBJ}
      & Top-$k$   & $O(n^2)$ & $1.00$ & $1.00$ & $1.00$
      & $1.00$ & $1.00$ & $1.00$
      & $0.48{\pm}0.06$ & $0.06{\pm}0.03$ & $0.00$ \\
      & Rerun     & $O(n^3)$ & $1.00$ & $1.00$ & $1.00$
      & $1.00$ & $1.00$ & $1.00$
      & $0.54{\pm}0.08$ & $0.06{\pm}0.02$ & $0.00$ \\
      & LPR       & $O(n^2)$ & $1.00$ & $1.00$ & $1.00$
      & $1.00$ & $1.00$ & $1.00$
      & $0.50{\pm}0.05$ & $0.05{\pm}0.05$ & $0.00$ \\
      & U-LPR     & $O(n^2)$ & $0.12{\pm}0.00$ & $0.00{\pm}0.01$ & $0.00$
      & $0.18{\pm}0.02$ & $0.00$ & $0.00$
      & $0.42{\pm}0.04$ & $0.00$ & $0.00$ \\
      \midrule
      \multirow{4}{*}{SGNN+U + SIT}
      & Top-$k$   & $O(n^2)$ & $1.00$ & $1.00$ & $1.00$
      & $0.55{\pm}0.40$ & $0.48{\pm}0.47$ & $0.33{\pm}0.58$
      & $0.34{\pm}0.11$ & $0.04{\pm}0.01$ & $0.00$ \\
      & Rerun     & $O(n^3)$ & $1.00$ & $1.00$ & $1.00$
      & $1.00$ & $1.00$ & $1.00$
      & $\mathbf{0.82{\pm}0.31}$ & $\mathbf{0.67{\pm}0.58}$ & $\mathbf{0.67{\pm}0.58}$ \\
      & LPR       & $O(n^2)$ & $1.00$ & $1.00$ & $1.00$
      & $0.55{\pm}0.40$ & $0.48{\pm}0.47$ & $0.33{\pm}0.58$
      & $0.34{\pm}0.11$ & $0.04{\pm}0.01$ & $0.00$ \\
      & U-LPR     & $O(n^2)$ & $0.11{\pm}0.01$ & $0.00$ & $0.00$
      & $0.18{\pm}0.03$ & $0.00$ & $0.00$
      & $0.29{\pm}0.04$ & $0.00$ & $0.00$ \\
      \midrule
      \multirow{4}{*}{SGNN+GU + SIT}
      & Top-$k$   & $O(n^2)$ & $1.00$ & $1.00$ & $1.00$
      & $1.00$ & $1.00$ & $1.00$
      & $0.47{\pm}0.05$ & $0.02{\pm}0.03$ & $0.00$ \\
      & Rerun     & $O(n^3)$ & $1.00$ & $1.00$ & $1.00$
      & $1.00$ & $1.00$ & $1.00$
      & $0.50{\pm}0.07$ & $0.03{\pm}0.03$ & $0.00$ \\
      & LPR       & $O(n^2)$ & $1.00$ & $1.00$ & $1.00$
      & $0.46{\pm}0.21$ & $0.37{\pm}0.33$ & $0.00$
      & $0.47{\pm}0.05$ & $0.02{\pm}0.03$ & $0.00$ \\
      & U-LPR     & $O(n^2)$ & $\mathbf{1.00}$ & $\mathbf{1.00}$ & $\mathbf{1.00}$
      & $\mathbf{1.00}$ & $\mathbf{1.00}$ & $\mathbf{1.00}$
      & $0.48{\pm}0.09$ & $0.04{\pm}0.04$ & $0.00$ \\
      \bottomrule
    \end{tabular}%
  }
\end{table}

\subsubsection{Rerun Decoding Is Strong but Expensive}

Rerun decoding gives perfect Easy and Medium recovery for every model, confirming
that a strong GNN ranking suffices when fresh scores are available after every
deletion~\citep{shoham2026ltr}. Its cost is consistently higher:
Table~\ref{tab:timing} reports $0.76$--$0.78$\,s per instance for Rerun, compared
with $0.20$--$0.21$\,s for LPR and $0.37$--$0.57$\,s for U-LPR.

\subsubsection{Without the Gradient Channel, U-LPR Collapses}

SGNN+U can still rank vertices reasonably under Top-$k$ and LPR, but its updater
does not work as a residual decoder. U-LPR stays near zero exact recovery: on Easy
it reaches only $0.11$--$0.13{\pm}0.01$ under different training regimes, while on
Medium it remains at $0.00$. SIT supervision alone is therefore not enough without
the gradient channel, and can even hurt U-LPR.

\subsubsection{The Gradient Channel Is Necessary and SIT Makes It Usable}

SGNN+GU under objective-only training still performs poorly under U-LPR
($0.12$--$0.18{\pm}0.02$ on Easy/Medium), showing that the gradient channel alone
is not enough---the model must also be trained to use it via SIT. Adding SIT
transforms U-LPR decoding to perfect approximation, overlap, and exact recovery on
Easy and Medium.

\subsubsection{SIT Trades Static Ranking for Update Quality}

For SGNN+GU, SIT hurts LPR ($1.00\to0.46$ on Medium) while enabling perfect U-LPR
recovery ($1.00$). The model sacrifices one-pass score quality to learn an effective
incremental update. For SGNN+U this trade-off does not occur: without the gradient
channel, SIT cannot teach $U_\theta$ a useful update rule, so both LPR and U-LPR
remain poor.

\subsubsection{Hard Regime}

Hard instances remain difficult for all models, consistent with the planted clique
hardness conjecture~\citep{manurangsi2021}. The strongest Hard result overall is
SGNN+U+SIT with Rerun decoding ($0.82{\pm}0.31$ approximation, $O(n^3)$). Among
$O(n^2)$ decoders, SGNN+GU+SIT with U-LPR is the strongest incremental variant
at $0.48{\pm}0.09$.

\begin{table}[H]
  \caption{Wall-clock time (seconds per instance) on planted-clique graphs with
    $n=1000$, mean $\pm$ std over 1000 held-out instances per regime, averaged over
  Easy, Medium, and Hard.}
  \label{tab:timing}
  \centering
  \begin{tabular}{lcccc}
    \toprule
    Model & Top-$k$ & Rerun & LPR & U-LPR \\
    \midrule
    SGNN + OBJ      & 0.06 & 0.78 & 0.21 & --- \\
    SGNN+U + OBJ    & 0.06 & 0.77 & 0.20 & 0.38 \\
    SGNN+GU + OBJ   & 0.06 & 0.78 & 0.21 & 0.57 \\
    SGNN+U + SIT    & 0.06 & 0.76 & 0.20 & 0.37 \\
    SGNN+GU + SIT   & 0.06 & 0.76 & 0.20 & 0.54 \\
    \bottomrule
  \end{tabular}
\end{table}

% ================================================================
\section{Interpretability and Concepts}
\label{sec:concepts}
% ================================================================

Table~\ref{tab:pca} reports PCA diagnostics, Spearman correlations, and the
above-median deletion concept diagnostic.

\begin{table}[H]
  \caption{Representation diagnostics on 1000 held-out test instances per regime
    ($n=1000$), reported as $|\rho|\pm\sigma$ (mean absolute Spearman $\pm$ std,
    all $p=0$ to machine precision unless noted).
    PC1/PC2: variance explained by the first two principal components of $h_v^{(L)}$.
    $|\rho|(\tilde{s},\text{deg})$: absolute Spearman between post-update scores and
    residual degree after one application of $U_\theta$ (E$=$Easy, M$=$Medium,
    H$=$Hard). \emph{Above-p15}: fraction of U-LPR steps where the removed vertex
    has residual degree above the 15th percentile among alive nodes (SIT training;
  medium regime; $10$ graphs $\times$ 500 steps, pooled). SGNN has no updater.}
  \label{tab:pca}
  \resizebox{\textwidth}{!}{%
    \begin{tabular}{llcccccccccccccc}
      \toprule
      & & \multicolumn{2}{c}{PC var.} &
      \multicolumn{3}{c}{$|\rho|(\text{PC1},\text{deg})$} &
      \multicolumn{3}{c}{$|\rho|(\text{logits},\text{deg})$} &
      \multicolumn{3}{c}{$|\rho|(\tilde{s},\text{deg})$} &
      \multicolumn{2}{c}{Above-p15 (\%)} \\
      \cmidrule(lr){3-4}\cmidrule(lr){5-7}\cmidrule(lr){8-10}\cmidrule(lr){11-13}\cmidrule(lr){14-15}
      Model & Train & PC1 & PC2 & E & M & H & E & M & H & E & M & H
      & Pooled & Mean$\pm$std \\
      \midrule
      SGNN & OBJ
      & 1.00 & 0.00
      & $1.00\pm0.00$ & $1.00\pm0.00$ & $1.00\pm0.00$
      & $1.00\pm0.00$ & $1.00\pm0.00$ & $1.00\pm0.00$
      & \multicolumn{3}{c}{--- (no updater)}
      & \multicolumn{2}{c}{---} \\
      \midrule
      \multirow{2}{*}{SGNN+U}
      & OBJ
      & 1.00 & 0.00
      & $1.00\pm0.00$ & $1.00\pm0.00$ & $1.00\pm0.00$
      & $1.00\pm0.00$ & $1.00\pm0.00$ & $1.00\pm0.00$
      & $0.52\pm0.01$ & $0.52\pm0.01$ & $0.50\pm0.01$
      & --- & --- \\
      & SIT
      & 1.00 & 0.00
      & $1.00\pm0.00$ & $1.00\pm0.00$ & $1.00\pm0.00$
      & $1.00\pm0.00$ & $1.00\pm0.00$ & $1.00\pm0.00$
      & $0.50\pm0.04$ & $0.50\pm0.03$ & $0.49\pm0.02$
      & $85.3$ & $85.3\pm1.3$ \\
      \midrule
      \multirow{2}{*}{SGNN+GU}
      & OBJ
      & 1.00 & 0.00
      & $1.00\pm0.00$ & $1.00\pm0.00$ & $1.00\pm0.00$
      & $1.00\pm0.00$ & $1.00\pm0.00$ & $1.00\pm0.00$
      & $0.51\pm0.01$ & $0.47\pm0.03$ & $0.36\pm0.17$
      & --- & --- \\
      & SIT
      & 1.00 & 0.00
      & $1.00\pm0.00$ & $1.00\pm0.00$ & $1.00\pm0.00$
      & $1.00\pm0.00$ & $1.00\pm0.00$ & $1.00\pm0.00$
      & $0.49\pm0.03$ & $0.49\pm0.04$ & $0.51\pm0.02$
      & $\mathbf{0.0}$ & $\mathbf{0.0\pm0.0}$ \\
      \bottomrule
    \end{tabular}%
  }
\end{table}

\subsection{The GNN Output Is Perfectly Degree-Ranked}

The final GNN output logits have absolute Spearman correlation $1.00{\pm}0.00$ with
vertex degree across all models and regimes, and PC1 explains essentially all
embedding variance (PC1$=1.00$, PC2$=0.00$), matching the degree-collapse finding
of~\citet{shoham2026ltr}. Figure~\ref{fig:pca_degree_embeddings} shows this geometry
for SGNN+U and SGNN+GU.

\begin{figure}[H]
  \centering
  \begin{minipage}{0.48\textwidth}
    \centering
    \includegraphics[width=\linewidth]{figures/SGNN_U_self_predictor_lpr_train_medium_g0_pc1_pc2_degree.pdf}
    \smallskip
    \textbf{(a)} SGNN+U.
  \end{minipage}\hfill
  \begin{minipage}{0.48\textwidth}
    \centering
    \includegraphics[width=\linewidth]{figures/SGNN_GU_self_predictor_lpr_train_medium_g0_pc1_pc2_degree.pdf}
    \smallskip
    \textbf{(b)} SGNN+GU.
  \end{minipage}
  \caption{Final-layer representation geometry. (\textbf{a}) SGNN+U and
    (\textbf{b}) SGNN+GU: vertices plotted in the first two principal-component
    coordinates of the last GNN layer, color-coded by residual degree. Both models
  show a strong degree--PC1 alignment before the updater is applied.}
  \label{fig:pca_degree_embeddings}
\end{figure}

\subsection{The Updater Learns a Low-Degree Pruning Concept}

After one application of $U_\theta$, the absolute Spearman correlation between
residual degree and updated scores drops to $0.36$--$0.52$, so the successful
updater is not merely preserving the raw degree ranking. SGNN+U+SIT removes an
above-15th-percentile-degree vertex on $85.3\%$ of U-LPR steps, while SGNN+GU+SIT
removes one on $0.0\%$ of steps. The learned concept is therefore a coarse
low-degree predicate: avoid deleting vertices above the 15th percentile of residual
degree among alive nodes. Figure~\ref{fig:score_degree_pruning_grid} shows the
pruning decisions at iterations $t=0,5,300$.
% Omitted for check: $t=1,10$.

\begin{figure}[H]
  \centering
  \begin{adjustbox}{max width=\textwidth}
    \begin{tabular}{@{}c c c c@{}}
      & $t=0$ & $t=5$ & $t=300$ \\
      % Omitted for check: & $t=1$ & $t=10$
      \rotatebox{90}{SGNN+U}
      & \includegraphics[width=0.30\textwidth]{figures/SGNN_U_self_predictor_lpr_train_medium_g0_upr_step0_score_vs_degree.pdf}
      & \includegraphics[width=0.30\textwidth]{figures/SGNN_U_self_predictor_lpr_train_medium_g0_upr_step5_score_vs_degree.pdf}
      & \includegraphics[width=0.30\textwidth]{figures/SGNN_U_self_predictor_lpr_train_medium_g0_upr_step300_score_vs_degree.pdf} \\
      % Omitted for check: & \includegraphics[width=0.30\textwidth]{figures/SGNN_U_self_predictor_lpr_train_medium_g0_upr_step1_score_vs_degree.pdf}
      % Omitted for check: & \includegraphics[width=0.30\textwidth]{figures/SGNN_U_self_predictor_lpr_train_medium_g0_upr_step10_score_vs_degree.pdf}
      \rotatebox{90}{SGNN+GU}
      & \includegraphics[width=0.30\textwidth]{figures/SGNN_GU_self_predictor_lpr_train_medium_g0_upr_step0_score_vs_degree.pdf}
      & \includegraphics[width=0.30\textwidth]{figures/SGNN_GU_self_predictor_lpr_train_medium_g0_upr_step5_score_vs_degree.pdf}
      & \includegraphics[width=0.30\textwidth]{figures/SGNN_GU_self_predictor_lpr_train_medium_g0_upr_step300_score_vs_degree.pdf}
      % Omitted for check: & \includegraphics[width=0.30\textwidth]{figures/SGNN_GU_self_predictor_lpr_train_medium_g0_upr_step1_score_vs_degree.pdf}
      % Omitted for check: & \includegraphics[width=0.30\textwidth]{figures/SGNN_GU_self_predictor_lpr_train_medium_g0_upr_step10_score_vs_degree.pdf}
    \end{tabular}
  \end{adjustbox}
  \caption{Score--degree pruning decisions during U-LPR. Rows correspond to
    SGNN+U and SGNN+GU; columns correspond to pruning iterations
    $t=0,5,300$. Each panel plots updater score (horizontal) vs.\ residual
    degree (vertical). $\times$ marks the vertex selected for removal; dashed line
  marks the 15th-percentile degree threshold.}
  \label{fig:score_degree_pruning_grid}
\end{figure}

% ================================================================
\section{Discussion}
\label{sec:discussion}
% ================================================================

The experiments point to a distinction that is easy to obscure when evaluating
learned combinatorial decoders: learning a good initial ranking is not the same as
learning a ranking that remains valid under residual updates. In our setting, the
GNN backbone reliably learns a degree-like ordering, and rerunning the GNN after
each deletion recovers strong Easy and Medium performance. Yet this does not by
itself yield an efficient decoder. The difficulty is not that the initial score is
uninformative; it is that the score is not intrinsically compatible with the residual
operation applied by the decoder.

This explains why the updater ablations are decisive. Without the gradient channel,
SIT gives the updater the right target sequence but not the right state variable.
The model is asked to imitate a rerun-based decoder from a nonlinear score whose
residual change is not directly available. This is why SGNN+U fails under U-LPR
even when trained with SIT. In contrast, SGNN+GU exposes a live gradient-derived
coordinate whose deletion effect is local and maintainable. SIT then has a usable
object to learn from: it does not need to invent the residual correction from a
static nonlinear score, but only to learn how to use the maintained gradient signal
to match the rerun decisions.

The results also clarify the role of cached maintenance. The cache is not merely an
implementation trick for speeding up the same computation. It changes what
information is available to the decoder throughout the pruning trajectory. A stale
gradient loses alignment because it ignores the vertices already removed; a
deletion-updated cache preserves the dominant residual change at $O(n)$ cost per
step. Thus the $O(n^2)$ decoder is not obtained by approximating an $O(n^3)$ rerun
more cheaply at each step, but by maintaining the right residual quantity
incrementally.

More broadly, the paper suggests that neural decoders should be judged not only by
score quality but by update compatibility. Classical residual algorithms are
efficient because their maintained statistics have explicit local update rules. Our
results show that a learned decoder can recover this property by exposing and
maintaining a gradient coordinate. The problem-specific part is the form of the
objective gradient and its cached sufficient statistics; the general principle is to
align the learned decoder with the local objective change caused by each deletion.

\subsection{Limitations and Future Work}

Our analysis suggests that gradient injection makes the residual correction
representable in the first linear layer of $U_\theta$, but it does not guarantee
that SGD will discover and use this consistently. A stronger solution would be an
architecture or loss that makes the GNN itself learn update-compatible coordinates,
ideally removing the need for a separate updater. Our training procedure also relies
on SIT supervision from a rerun-based teacher ($T=300$ removals per instance), which
is expensive even though inference is quadratic. Finally, the experiments focus on
planted clique; extending the gradient-maintenance principle to other locally
decomposable objectives such as independent set and constraint-satisfaction
relaxations remains an important direction.

% ================================================================
%  MDPI REQUIRED END-MATTER
% ================================================================

\authorcontributions{%
  Conceptualization, E.S. and D.V.; methodology, E.S. and D.V.; software, E.S.;
  validation, E.S., H.R. and D.V.; formal analysis, E.S. and D.V.; investigation,
  E.S.; writing---original draft preparation, E.S.; writing---review and editing,
  H.R. and D.V.; visualization, E.S.; supervision, D.V.
  All authors have read and agreed to the published version of the manuscript.
}

\funding{This research received no external funding.}

\institutionalreview{Not applicable.}

\informedconsent{Not applicable.}

\dataavailability{%
  The code and data supporting the reported results will be made available at
  \url{https://github.com/TODO} upon publication.
}


\conflictsofinterest{The authors declare no conflicts of interest.}

% ================================================================
%  ABBREVIATIONS
% ================================================================
\abbreviations{Abbreviations}{%
  The following abbreviations are used in this manuscript:\\
  \noindent
  \begin{tabular}{@{}ll}
    GNN     & Graph Neural Network \\
    MPNN    & Message-Passing Neural Network \\
    MLP     & Multi-Layer Perceptron \\
    LPR     & Least Probable Removal \\
    LDR     & Least Degree Removal \\
    U-LPR   & Updatable Least Probable Removal \\
    SIT     & Self-Imitation Training \\
    OBJ     & Objective (training regime) \\
    SGNN    & Scattering GNN (backbone only) \\
    SGNN+U  & SGNN with score updater, no gradient channel \\
    SGNN+GU & SGNN with score updater and gradient channel \\
    WL      & Weisfeiler--Leman (graph isomorphism test) \\
    SDP     & Semidefinite Program \\
    UGC     & Unique Games Conjecture \\
    PC1/PC2 & First/Second Principal Component \\
  \end{tabular}
}

% ================================================================
\bibliography{bib}
% (mdpi.cls sets \bibliographystyle{Definitions/mdpi} automatically)

% ================================================================
\clearpage
\appendixtitles{yes}
\appendix
\section{Algorithms}
\label{app:algorithms}
% ================================================================

\begin{algorithm}[H]
  \caption{U-LPR: Cached-Deletion Decoder}
  \label{alg:gub}
  \begin{algorithmic}[1]
    \State Run GNN once to obtain scores $s_i$ and probabilities $p_i$ for all $i$.
    \State Initialize gradient/cache from $p$. \hfill\Comment{$O(n^2)$}
    \State $C \leftarrow V(G)$;\; $H \leftarrow []$.
    \While{$C$ does not form a clique}
    \State $v^* \leftarrow \arg\min_{i \in C} s_i$;\; append $v^*$ to $H$;\;
    $C \leftarrow C \setminus \{v^*\}$.
    \State Remove $v^*$'s contribution from gradient cache. \hfill\Comment{$O(n)$}
    \State $g_i \leftarrow$ gradient from updated cache for all $i \in C$.
    \hfill\Comment{$O(n)$}
    \State $s_i \mathrel{+}= U_\theta(s_i,\, s_{v^*},\, A_{iv^*},\, g_i,\,
    g_{v^*})$ for all $i \in C$. \hfill\Comment{$O(n)$}
    \EndWhile
    \For{$v$ from last to first in $H$}
    \If{$C \cup \{v\}$ is a clique} $C \leftarrow C \cup \{v\}$. \EndIf
    \EndFor
    \State \Return $C$.
  \end{algorithmic}
\end{algorithm}

\begin{algorithm}[H]
  \caption{SIT Training Step}
  \label{alg:sit_training}
  \begin{algorithmic}[1]
    \State Sample instance $(G, k)$; run GNN once to obtain $\tilde{s}$, $p$.
    \State Initialize gradient/cache from $p$.
    \hfill\Comment{$O(n^2)$; SGNN+GU only}
    \For{$t = 0, \ldots, T-1$}
    \State $v^t_{\mathrm{teach}} \leftarrow \arg\min_{i \in V(G_t)}
    \mathrm{GNN}(G_t)_i$. \hfill\Comment{Teacher reruns GNN, $O(n^2)$}
    \State $\mathcal{L}_{\mathrm{SIT}} \mathrel{+}=
    \mathrm{CE}_{\mathrm{softmax}}(-\tilde{s}_{V(G_t)},\,
    v^t_{\mathrm{teach}})$.
    \hfill\Comment{Student predicts same removal}
    \State $G_t \leftarrow G_t \setminus \{v^t_{\mathrm{teach}}\}$.
    \State Update gradient cache by removing $v^t_{\mathrm{teach}}$.
    \hfill\Comment{$O(n)$; SGNN+GU only}
    \State $\tilde{s}_i \mathrel{+}= U_\theta(\tilde{s}_i,\,
    \tilde{s}_{v^t},\, A_{iv^t},\, \gamma_i,\, \gamma_{v^t})$
    for all $i \in V(G_t)$.
    \EndFor
    \State Optimise $\mathcal{L}_{\mathrm{clique}} +
    \mathcal{L}_{\mathrm{SIT}} / T$.
  \end{algorithmic}
\end{algorithm}

\end{document}
